% =====================================================================
% Judging_by_the_Cover_v35_section5_patch.tex
% =====================================================================
% Drop-in patches against the v34 Overleaf source. The source file is
% on Overleaf (not in this repo) and the in-session paste was truncated
% mid-Section 5, so this file is a patch rather than a full manuscript.
% Apply each labeled block in Overleaf.
%
% Scope (per the §5 rewrite spec):
%   [A] Abstract -- drop-in replacement of the "three classifier
%       families" block, keeping paraphrase-infeasibility sentence.
%   [B] Introduction contributions item (3) -- drop-in replacement.
%   [C] Section 5 header -- remove the n=800 TODO.
%   [D] Section 5.1 -- UNCHANGED (v34 singleton-probe subsection).
%   [E] NEW Section 5.2 "Bilateral singleton probe with length matching"
%       incl. new Table~\ref{tab:bilateral_probe}.
%   [F] Section 5.3 = old Section 5.2 (paraphrase infeasibility),
%       renumbered only; content unchanged.
%   [G] Conclusion -- sentence-level edits.
%   [H] Limitations -- sentence-level additions.
%   [I] Diff summary.
%
% Rules observed: no retraining, no API calls, no changes to §3/§4/§6
% or any table therein, exploratory framing for the scale pattern,
% same-index as the lead metric.
% =====================================================================


% =====================================================================
% [A] ABSTRACT -- drop-in replacement
% Replace the block that currently begins
%   "We further probe classifier robustness by training each of three
%   classifier families..."
% through the end of the paraphrase-infeasibility sentence.
% =====================================================================

% >>> BEGIN [A] <<<
We further probe classifier robustness with three adversarial probes.
A singleton probe against $20$ adversarially generated false
statements with truthful-side surface features shows reduced
susceptibility for classifiers trained on the cleaned subset across
three families (interpretable surface-feature, BGE-large,
ModernBERT-base), with ModernBERT-base exhibiting the largest mean
reduction on that probe ($+0.135$). A bilateral singleton probe with
length-matched true/false pairs extends the audit to seven classifier
families spanning $149$M to $3.8$B parameters (the interpretable
$10$-feature model plus BGE-large, ModernBERT-base, Qwen2.5-0.5B,
SmolLM2-1.7B, Qwen2.5-3B, and Phi-3.5-mini); on this probe the
interpretable surface classifier shows the largest cleaning benefit
(bilateral fool-drop $+0.35$ same-index, $+0.45$ all-pairs at $n=20$),
and two $3$B+ causal-LM families show substantially larger fool-drops
than sub-$2$B classifiers ($+0.30$ vs.\ $+0.05$--$+0.10$ same-index),
suggesting a scale-dependent cleaning effect that warrants
larger-sample validation. A complementary paired-paraphrase probe
reveals that truth-preserving surface-form flip is structurally
infeasible on TruthfulQA, because the correct-side cues are natural
language operators with truth-conditional content rather than
stylistic ornaments.
% >>> END [A] <<<


% =====================================================================
% [B] INTRODUCTION -- contributions list item (3), drop-in replacement
% Replace the existing "(3) We construct an adversarial probe..."
% bullet in the contributions list.
% =====================================================================

% >>> BEGIN [B] <<<
(3) We construct adversarial probes of the surface-form shortcut at
three levels. A \emph{singleton probe} tests three classifier
families (interpretable surface-feature, BGE-large, ModernBERT-base)
against $20$ adversarially generated false statements with
truthful-side surface features; classifiers trained on the cleaned
subset show reduced susceptibility across all three families
(ModernBERT-base mean reduction $+0.135$). A \emph{bilateral
singleton probe with length matching} extends the audit to seven
classifier families spanning $149$M to $3.8$B parameters, pairing
$20$ length-matched true statements with the singleton-probe false
statements. The interpretable surface classifier shows the largest
bilateral cleaning benefit (same-index fool-drop $+0.35$, all-pairs
$+0.45$), sub-$2$B semantic encoders show consistent modest
fool-drops ($+0.05$ to $+0.10$ same-index), and two independent
$3$B+ causal-LM families (Qwen2.5-3B, Phi-3.5-mini) show
substantially larger same-index fool-drops ($+0.30$); we present
this scale pattern as exploratory given the $n=20$ pair budget. A
\emph{paired paraphrase probe} reveals that truth-preserving
surface-form flip is structurally infeasible on TruthfulQA, because
the correct-side cues (negation, hedging, authority attribution) are
natural language operators with truth-conditional content rather
than stylistic ornaments.
% >>> END [B] <<<


% =====================================================================
% [C] SECTION 5 HEADER -- replace the n=800 TODO paragraph
% The v34 §5 preamble ends with a sentence stating
%   "A scale-up to n=800 is plan[ned]..."
% Replace that sentence with the honest n=20 framing below; keep the
% rest of the preamble unchanged.
% =====================================================================

% >>> BEGIN [C] <<<
All quantitative results in this section are based on a small Stage~0
sample of $n=20$ adversarial items per probe. This sample size is
sufficient to establish direction and rough effect size but admits
wide confidence intervals; we therefore treat every numerical
comparison in this section as suggestive rather than decisive, and
flag scale-dependent patterns as exploratory where they arise.
% >>> END [C] <<<


% =====================================================================
% [D] SECTION 5.1 -- UNCHANGED
% Keep the v34 §5.1 "Singleton probe..." subsection (including its
% existing 3-family table) exactly as-is.
% =====================================================================


% =====================================================================
% [E] NEW SECTION 5.2 -- full subsection
% Insert directly after the v34 §5.1 subsection, before the old §5.2
% (which becomes §5.3 in [F]).
% =====================================================================

% >>> BEGIN [E] <<<
\subsection{Bilateral singleton probe with length matching}
\label{sec:bilateral-probe}

The singleton probe above measures classifier vulnerability to a
single-sided set of adversarial false statements carrying
truthful-side surface cues. A complementary test is whether
classifiers trained on the audited subset can be \emph{fooled} when a
surface-plain true statement is presented alongside a false statement
dressed in truthful-side surface form. We construct such a bilateral
probe and evaluate it across an expanded ladder of seven classifier
families spanning four model sizes and three pretraining paradigms.

\paragraph{Setup.} We generate $20$ true statements (Claude Opus~$4.7$,
judged by GPT-$5.4$ pinned at \texttt{gpt-5.4-2026-03-05}; $19/20$
judged true) stripped of all truthful-side surface cues: no negation
openers or tokens, no hedging, no authority attribution, and no
contrast conjunctions. Each A-side true statement is length-matched to
the $\S$5.1 B-side false statements, with word count constrained to
the range $[22, 26]$ to match the B-side mean of $23.5$ words. The
$20$ A-side true statements are paired with the $20$ existing B-side
false statements from the singleton probe. For a classifier $f$, the
\emph{fool rate} on a pair set is the fraction of pairs in which $f$
assigns higher $P(\text{truthful})$ to the B-side (surface-true,
semantically false) than to the A-side (surface-plain, semantically
true). We report fool rates under two pairings: \emph{same-index}
($n=20$), which preserves topic alignment between $A_i$ and $B_i$ and
which we treat as the lead metric; and \emph{all-pairs} ($n=400$),
the $20 \times 20$ cross-product, which we treat as a secondary
robustness check because it entangles topic coherence with truth
discrimination. For each classifier, the \emph{fool drop} is the
difference $\text{fool}_\text{full} - \text{fool}_\text{cleaned}$;
positive values indicate that cleaning reduces bilateral shortcut
reliance.

\paragraph{Classifier ladder.} We evaluate seven families: the
interpretable $10$-feature surface classifier from $\S$3; BGE-large
($335$M, contrastive encoder); ModernBERT-base ($149$M, masked-LM
encoder); and four causal-LM families used as frozen mean-pooled
encoders---Qwen2.5-0.5B-Instruct ($494$M), SmolLM2-1.7B-Instruct
($1.7$B), Qwen2.5-3B-Instruct ($3.1$B, bf16), and Phi-3.5-mini-Instruct
($3.8$B, bf16). Every embedding-based family uses the same
\texttt{StandardScaler}${+}$\texttt{LogisticRegression} head protocol,
fit independently on the full $790$-pair TruthfulQA and the cleaned
$\tau=0.52$ $528$-pair subset, with grouped $5$-fold cross-validation
by \texttt{pair\_id}.

\begin{table}[H]
\centering
\small
\caption{Bilateral singleton probe across seven classifier families
($n=20$ same-index, $n=400$ all-pairs). CV AUC columns report
grouped $5$-fold cross-validation AUC on full and cleaned
($\tau=0.52$) TruthfulQA. Fool-drop $=$ full fool-rate $-$ cleaned
fool-rate; positive values indicate cleaning reduces bilateral
shortcut reliance. Same-index is the lead metric; all-pairs is
secondary and mixes topic coherence with truth discrimination.}
\label{tab:bilateral_probe}
\begin{tabular}{@{}lccc@{}}
\toprule
Classifier & CV AUC (full / cleaned) & Same-idx fool-drop & All-pairs fool-drop \\
\midrule
surface\_lr ($10$ feats)    & $0.718$ / $0.596$ & $+0.350$ & $+0.447$ \\
BGE-large ($335$M)          & $0.785$ / $0.720$ & $+0.050$ & $+0.088$ \\
ModernBERT-base ($149$M)    & $0.749$ / $0.711$ & $+0.100$ & $+0.057$ \\
Qwen2.5-0.5B ($494$M)       & $0.746$ / $0.717$ & $+0.100$ & $+0.087$ \\
SmolLM2-1.7B                & $0.795$ / $0.755$ & $+0.050$ & $+0.078$ \\
Qwen2.5-3B ($3.1$B)         & $0.810$ / $0.806$ & $+0.300$ & $+0.178$ \\
Phi-3.5-mini ($3.8$B)       & $0.805$ / $0.782$ & $+0.300$ & $+0.325$ \\
\bottomrule
\end{tabular}
\end{table}

\paragraph{Results.} Every family shows a non-negative same-index
fool-drop, indicating that cleaning reduces bilateral susceptibility
across all seven classifiers. The interpretable surface classifier
shows the largest effect ($+0.350$ same-index, $+0.447$ all-pairs),
consistent with its direct dependence on the lexical features
targeted by the $\tau=0.52$ audit. The four frozen sub-$2$B semantic
encoders (BGE-large, ModernBERT-base, Qwen2.5-0.5B, SmolLM2-1.7B)
show modest but consistent fool-drops of $+0.05$ to $+0.10$
same-index.

\paragraph{Scale pattern (exploratory).} At $3$B+ parameters,
Qwen2.5-3B and Phi-3.5-mini show substantially larger bilateral
fool-drops than the sub-$2$B classifiers ($+0.30$ vs.\ $+0.05$--$+0.10$
same-index). Whether this reflects a genuine scale threshold requires
larger-sample validation; the effect is replicated across two
independent causal-LM families at matched scale, but $n=20$ admits
wide confidence intervals, and the two families diverge on the
all-pairs metric (Phi-3.5-mini $+0.325$ vs.\ Qwen2.5-3B $+0.178$),
which may reflect Phi-specific synthetic-training-data amplification
rather than a pure scale effect. We therefore present this as a
suggestive observation rather than an established threshold.

\paragraph{Length matching as a methodological contribution.} The
v7a-v3 probe required constraining A-side word count to match the
B-side distribution ($22$--$26$ words, centered on the B-side mean of
$23.5$ words). At shorter A-side lengths ($5$--$15$ words), the
cleaned \texttt{surface\_lr} classifier's residual \texttt{word\_count}
feature---whose standardized coefficient rose from the $4$th-largest
in \texttt{surface\_lr\_full} ($+0.226$) to the $1$st-largest in
\texttt{surface\_lr\_cleaned} ($+0.311$), while \texttt{neg\_cnt}
dropped from $+0.746$ to $+0.045$---masked the cleaning effect on the
bilateral probe. Neutralizing this residual length shortcut by
length-matching A-sides to B-sides revealed the fool-drops reported
in Table~\ref{tab:bilateral_probe}. This pattern indicates that
single-cue cleaning procedures can transfer shortcut reliance from
the targeted feature to the next-available residual feature, and
that probes of cleaned classifiers should explicitly control for
candidate residual cues rather than treat the post-cleaning
feature distribution as stationary. We view this as a
methodological contribution of the bilateral probe: adversarial
probes of surface-cleaned classifiers should neutralize features that
were not the direct target of the cleaning procedure, or risk
conflating a residual shortcut with the absence of a cleaning effect.
% >>> END [E] <<<


% =====================================================================
% [F] OLD SECTION 5.2 --> SECTION 5.3 (renumber only)
% The v34 "paired paraphrase / infeasibility" subsection is kept
% verbatim, moved to follow the new §5.2 above. If you use an explicit
% \label, rename from sec:paraphrase-probe (or similar) to
% sec:paraphrase-probe unchanged -- just re-position in file order.
% No textual edits.
% =====================================================================


% =====================================================================
% [G] CONCLUSION -- sentence-level edits
% In the v34 conclusion, the existing adversarial-probing sentence
% refers to "three classifier families". Replace that sentence with
% [G1] and append [G2] + [G3] after it.
% =====================================================================

% >>> BEGIN [G1] (replaces the "three classifier families" sentence) <<<
Adversarial probing across three complementary designs---a
singleton probe (three classifier families), a bilateral
length-matched singleton probe (seven classifier families spanning
$149$M to $3.8$B parameters, across engineered features,
contrastive encoder, masked-LM encoder, and causal-LM paradigms),
and a paired-paraphrase probe---confirms that surface-form cleaning
reduces shortcut reliance across every family tested, with the
largest bilateral cleaning benefit on the interpretable lexical
classifier (same-index fool-drop $+0.35$) and smaller but
consistent benefits on frozen semantic encoders.
% >>> END [G1] <<<

% >>> BEGIN [G2] (new sentence, insert immediately after [G1]) <<<
Two independent causal-LM families at $3$B+ parameters show
substantially larger bilateral fool-drops than sub-$2$B classifiers,
but $n=20$ admits wide confidence intervals and we leave
scale-threshold claims for larger-sample work.
% >>> END [G2] <<<

% >>> BEGIN [G3] (new sentence, insert immediately after [G2]) <<<
The bilateral probe also surfaces a methodological point: the
cleaning procedure transferred \texttt{surface\_lr}'s dominant
coefficient from negation to word count, so credible adversarial
probes of cleaned classifiers must length-match---or otherwise
neutralize---features that were not the direct target of the
cleaning procedure.
% >>> END [G3] <<<


% =====================================================================
% [H] LIMITATIONS -- additions
% Keep the existing n=20 caveat. Append [H1] and [H2] to the
% adversarial-probing limitations paragraph.
% =====================================================================

% >>> BEGIN [H1] (new sentence) <<<
The $3$B+ scale pattern reported in $\S$5.2---two causal-LM
families showing substantially larger same-index fool-drops than the
five sub-$2$B classifiers---is based on $n=20$ pairs and on only two
families above the apparent threshold. We therefore present it as
exploratory and consistent with, but not evidence of, a
capability-dependent cleaning effect; larger samples and additional
$3$B+ families are required before interpreting this as a scale
threshold.
% >>> END [H1] <<<

% >>> BEGIN [H2] (new sentence) <<<
The length-matching requirement in $\S$5.2 arose from a
cleaning-induced coefficient shift in \texttt{surface\_lr} (negation
$\to$ word count). This is a property of the specific single-cue
$\tau=0.52$ audit used in this paper rather than a universal
probing procedure; future cleaning protocols may produce different
residual shortcut features, and bilateral probes should be tuned to
whichever residual cue dominates the cleaned classifier under study.
% >>> END [H2] <<<


% =====================================================================
% [J] PREREG AMENDMENT -- add 8th classifier family (Qwen2.5-1.5B)
% After freezing the v35 draft we added Qwen2.5-1.5B-Instruct as an 8th
% classifier family, to fill the scale gap between Qwen2.5-0.5B (494M)
% and SmolLM2-1.7B / Qwen2.5-3B in the bilateral probe. Same recipe as
% the other causal-LM families (mean-pooled frozen encoder, bf16 on MPS,
% StandardScaler + LR head, grouped 5-fold CV by pair_id).
%
% This block is reported in full regardless of outcome, per the
% prereg-amendment protocol for adding a family after the original
% table was frozen. The n=20 v7a-v3 bilateral probe in Table
% \ref{tab:bilateral_probe} is UNCHANGED; this amendment only adds the
% scaled v9 same-question results.
% =====================================================================

% >>> BEGIN [J1] (amendment: family ladder extended to eight) <<<
\paragraph{Amendment: eighth classifier family.} After the original
seven-family ladder was frozen, we added Qwen2.5-1.5B-Instruct
($1.5$B, causal LM, mean-pool, bf16) to fill the scale gap between
Qwen2.5-0.5B and SmolLM2-1.7B. The model is trained with the
identical \texttt{StandardScaler}${+}$\texttt{LogisticRegression}
recipe on the same full ($790$-pair) and cleaned ($\tau=0.52$,
$528$-pair) TruthfulQA splits, with grouped $5$-fold cross-validation
by \texttt{pair\_id}. Grouped-CV AUC: full $0.770$, cleaned $0.738$.
We report its bilateral probe result regardless of outcome in the
scaled evaluation below.
% >>> END [J1] <<<

% >>> BEGIN [J2] (new Table: v9 same-question scaled bilateral probe) <<<
\begin{table}[H]
\centering
\small
\caption{Scaled bilateral same-question probe across nine classifier
families ($n=206$ topic-matched true/false pairs; judge-verified on
both sides with \texttt{gpt-5.4-2026-03-05}).
Accuracy $=$ fraction of pairs where the classifier ranks the true
side above the false side; chance $=0.5$.
$\Delta = \text{Acc}_\text{cleaned} - \text{Acc}_\text{full}$;
positive values indicate cleaning reduces surface-shortcut reliance.
$p$ columns are two-sided exact binomial vs.\ $0.5$ (Acc columns) and
exact McNemar on pick-flips (cleaning column). Amendment families
Qwen2.5-1.5B marked $\dagger$ and Llama-3.2-3B marked $\ddagger$.}
\label{tab:bilateral_probe_scaled}
\begin{tabular}{@{}lccccc@{}}
\toprule
Classifier & Acc$_\text{full}$ [95\% CI] & Acc$_\text{cleaned}$ [95\% CI] & $\Delta$ & McNemar $p$ & flips $+/-$ \\
\midrule
surface\_lr ($10$ feats)    & $0.000$ [$0.00$, $0.02$] & $0.005$ [$0.00$, $0.03$] & $+0.005$ & $1.000$    & $+1/-0$  \\
BGE-large ($335$M)          & $0.451$ [$0.39$, $0.52$] & $0.495$ [$0.43$, $0.56$] & $+0.044$ & $0.222$    & $+26/-17$ \\
ModernBERT-base ($149$M)    & $0.587$ [$0.52$, $0.65$] & $0.549$ [$0.48$, $0.61$] & $-0.039$ & $0.322$    & $+21/-29$ \\
Qwen2.5-0.5B ($494$M)       & $0.549$ [$0.48$, $0.61$] & $0.549$ [$0.48$, $0.61$] & $+0.000$ & $1.000$    & $+21/-21$ \\
Qwen2.5-1.5B ($1.5$B)$^\dagger$ & $0.549$ [$0.48$, $0.61$] & $0.568$ [$0.50$, $0.63$] & $+0.019$ & $0.636$    & $+22/-18$ \\
SmolLM2-1.7B                & $0.505$ [$0.44$, $0.57$] & $0.578$ [$0.51$, $0.64$] & $+0.073$ & $0.024^{*}$  & $+27/-12$ \\
Qwen2.5-3B ($3.1$B)         & $0.500$ [$0.43$, $0.57$] & $0.597$ [$0.53$, $0.66$] & $+0.097$ & $0.006^{**}$ & $+34/-14$ \\
Llama-3.2-3B ($3.2$B)$^\ddagger$ & $0.437$ [$0.37$, $0.51$] & $0.519$ [$0.45$, $0.59$] & $+0.083$ & $0.006^{**}$ & $+26/-9$  \\
Phi-3.5-mini ($3.8$B)       & $0.578$ [$0.51$, $0.64$] & $0.583$ [$0.51$, $0.65$] & $+0.005$ & $1.000$    & $+20/-19$ \\
\bottomrule
\end{tabular}
\end{table}
% >>> END [J2] <<<

% >>> BEGIN [J3] (amendment result paragraph) <<<
\paragraph{Amendment result.} Qwen2.5-1.5B shows a small positive
bilateral $\Delta$ of $+0.019$ (Acc$_\text{full}=0.549$,
Acc$_\text{cleaned}=0.568$), with McNemar $p=0.636$; the effect is
not distinguishable from zero at the scaled $n=206$ sample. The 1.5B
family falls in the same "no-big-effect" plateau as Qwen2.5-0.5B
($+0.000$) and Phi-3.5-mini ($+0.005$), which tightens the
interpretation of the scale pattern from $\S$5.2: at the scaled
sample size, only Qwen2.5-3B and SmolLM2-1.7B retain significant
cleaning benefit (McNemar $p=0.006$ and $p=0.024$ respectively), and
the original "3B+ threshold" framing is not supported---Qwen2.5-1.5B
and Phi-3.5-mini behave alike despite the $2.5{\times}$ size
difference. We report the null result in full and treat the
$n=20$ scale narrative from $\S$5.2 as superseded by the
$n=206$ evaluation in Table~\ref{tab:bilateral_probe_scaled}.
% >>> END [J3] <<<


% =====================================================================
% [K] PREREG AMENDMENT -- add 9th classifier family (Llama-3.2-3B)
% After the eighth-family amendment above, we added
% meta-llama/Llama-3.2-3B-Instruct as an architecture-matched replicate
% of Qwen2.5-3B at near-identical scale (3.2B vs 3.1B). The scientific
% purpose is to test whether the Qwen2.5-3B cleaning-helps pattern in
% Table \ref{tab:bilateral_probe_scaled} is specific to the Qwen
% architecture or replicates in a different model family at the same
% scale. Same recipe as the other causal-LM families (mean-pooled
% frozen encoder, bf16 on MPS, StandardScaler + LR head, grouped
% 5-fold CV by pair_id). Reported in full regardless of outcome.
% The n=20 v7a-v3 bilateral probe in Table \ref{tab:bilateral_probe}
% is UNCHANGED. Table \ref{tab:bilateral_probe_scaled} has been
% extended with a ninth row for Llama-3.2-3B (marked \ddagger) under
% the same caption/label.
% =====================================================================

% >>> BEGIN [K1] (amendment: family ladder extended to nine) <<<
\paragraph{Amendment: ninth classifier family.} Following the
eighth-family amendment above, we added
\texttt{meta-llama/Llama-3.2-3B-Instruct} ($3.2$B, causal LM,
mean-pool, bf16) as an architecture-matched replicate of Qwen2.5-3B
at near-identical scale. The scientific purpose is to test whether
the Qwen2.5-3B cleaning effect observed in
Table~\ref{tab:bilateral_probe_scaled} is specific to the Qwen
architecture or replicates in a different model family at the same
$3$B scale. Recipe, splits, and CV grouping are identical to the
other causal-LM families. Grouped-CV AUC: full $0.812$, cleaned
$0.780$---the highest full-data CV AUC of all nine families. We
report its bilateral probe result regardless of outcome.
% >>> END [K1] <<<

% >>> BEGIN [K3] (amendment result paragraph for Llama-3.2-3B) <<<
\paragraph{Amendment result.} Llama-3.2-3B shows
$\text{Acc}_\text{full}=0.437$ [$0.37$, $0.51$] (near-significantly
below chance, two-sided binomial $p=0.081$) and
$\text{Acc}_\text{cleaned}=0.519$ [$0.45$, $0.59$], yielding
$\Delta=+0.083$ with exact McNemar $p=0.006^{**}$ (pick-flips
$+26/-9$). This replicates the Qwen2.5-3B pattern---same-scale
replicate, $\Delta=+0.097$, McNemar $p=0.006^{**}$---in a different
model family. With this addition, three of the nine families in
Table~\ref{tab:bilateral_probe_scaled} show exact-McNemar-significant
cleaning benefit at $n=206$: Qwen2.5-3B, Llama-3.2-3B (both
$3$B-class causal LMs from different families), and SmolLM2-1.7B.
Two same-scale $3$B-class models from unrelated pre-training lineages
(Qwen, Llama) thus produce concordant cleaning effects, while
neighbouring scales (1.5B, 3.8B) and smaller causal-LM families
(0.5B) do not, indicating a model-family-independent effect
localised near the $3$B scale rather than a monotone scale trend.
We retain the $n=206$ evaluation in
Table~\ref{tab:bilateral_probe_scaled} as the primary scaled result.
% >>> END [K3] <<<


% =====================================================================
% [I] DIFF SUMMARY
% =====================================================================
%
%  Section / block                       Action
%  ------------------------------------- ------------------------------
%  Abstract "three classifier families"  Replace with [A]
%  Intro contributions item (3)          Replace with [B]
%  §5 preamble (n=800 TODO sentence)     Replace with [C]
%  §5.1 singleton probe                  UNCHANGED
%  §5.2 NEW bilateral probe              Insert [E]
%    - new Table \ref{tab:bilateral_probe}  (7-family, 3 columns)
%    - 4 paragraphs (setup, ladder, results, scale, length-matching)
%  §5.2 old (paraphrase infeasibility)   Renumber to §5.3; no text edits
%  Conclusion "three classifier families" Replace with [G1]
%  Conclusion (appended)                 Add [G2] and [G3]
%  Limitations (appended)                Add [H1] and [H2]
%  Prereg amendment (new §5.2.1)          Insert [J1] + [J2] + [J3]
%    - new Table \ref{tab:bilateral_probe_scaled}  (9-family, n=206)
%  Prereg amendment (new §5.2.2)          Insert [K1] + [K3]
%    - extends Table \ref{tab:bilateral_probe_scaled} with Llama-3.2-3B
%  §3, §4, §6, all existing tables       UNCHANGED
%
% Tables added: 2 (tab:bilateral_probe, tab:bilateral_probe_scaled)
% Tables changed: 0
% Tables removed: 0
% Figures added/changed/removed: 0
% =====================================================================
